\GalSourcePageBoundary{069}{L0068}{40}{40}

puisque, dans le théorème précédent, il ne s'agit que des substitutions
que l'on peut faire dans les fonctions.

\emph{Scolie II.} --- Les substitutions sont indépendantes même du
nombre des racines.

\begin{center}{\bfseries PROPOSITION II.}\end{center}
\GalSourceErrorMark{W09-SE002}

\noindent\textsc{Théorème}\textsuperscript{(*)}. --- \emph{Si l'on adjoint à une équation donnée la
racine $r$ d'une équation auxiliaire irréductible, 1\textsuperscript{o} il arrivera
de deux choses l'une: ou bien le groupe de l'équation ne sera
pas changé, ou bien il se partagera en $p$ groupes appartenant
chacun à l'équation proposée respectivement quand on lui adjoint
chacune des racines de l'équation auxiliaire; 2\textsuperscript{o} ces
groupes jouiront de la propriété remarquable, que l'on passera
de l'un à l'autre en opérant dans toutes les permutations du
premier une même substitution de lettres.}

1\textsuperscript{o} Si, après l'adjonction de $r$, l'équation en $V$, dont il est
question plus haut, reste irréductible, il est clair que le groupe
de l'équation ne sera pas changé. Si, au contraire, elle se réduit,
alors l'équation en $V$ se décomposera en $p$ facteurs, tous de même
degré et de la forme
\[
f(V,r)\times f(V,r')\times f(V,r'')\times\ldots,
\]
$r,r',r'',\ldots$ étant d'autres valeurs de $r$. Ainsi, le groupe de
l'équation proposée se décomposera aussi en groupes chacun
d'un même nombre de permutations, puisque à chaque valeur
de $V$ correspond une permutation. Ces groupes seront respecti-

\GalHistoricalNote{\textsuperscript{(*)} Dans l'énoncé du théorème, après ces mots: \emph{la racine $r$ d'une équation
auxiliaire irréductible}, Galois avait mis d'abord ceux-ci: \emph{de degré $p$ premier},
qu'il a effacés plus tard. De même, dans la démonstration, au lieu de $r,r',r'',\ldots$
étant d'autres valeurs de $r$, la rédaction primitive portait: $r,r',r'',\ldots$
étant les diverses valeurs de $r$. Enfin on trouve à la marge du manuscrit la
note suivante de l'auteur:

« Il y a quelque chose à compléter dans cette démonstration. Je n'ai pas le
temps. »

Cette ligne a été jetée avec une grande rapidité sur le papier; circonstance
qui, jointe aux mots: « Je n'ai pas le temps », me fait penser que Galois a relu
son Mémoire pour le corriger avant d'aller sur le terrain. \hfill \textsc{(A.~Ch.)}}
